\documentclass{article}

% Options:
%   [final]         - Non-anonymous submission showing student names
%   [nonanonymous]  - Explicitly non-anonymous
%   (default)       - Anonymous (useful if doing blind peer-grading)
\usepackage[final]{deepgate_report}

\usepackage[utf8]{inputenc}
\usepackage[T1]{fontenc}
\usepackage{hyperref}
\usepackage{url}
\usepackage{booktabs}
\usepackage{amsfonts}
\usepackage{amsmath}
\usepackage{nicefrac}
\usepackage{microtype}
\usepackage{graphicx}
\usepackage{subcaption}

\title{DeepGate-FaceRec: Open-Set Face Verification via Deep Metric Learning and Online Triplet Mining}

\author{%
  Firstname Lastname \\
  Student ID: 12345678 \\
  Department of Computer Science \\
  \texttt{student@university.edu} \\
  \And
  Partner Firstname Lastname \\
  Student ID: 87654321 \\
  Department of Computer Science \\
  \texttt{partner@university.edu} \\
}

\begin{document}

\maketitle

\begin{abstract}
Traditional softmax classification heads are constrained to closed-set problem formulations, rendering them ineffective for biometric access systems where enrolled identities dynamically expand without retraining. In this report, we implement and evaluate a deep metric learning framework for open-set face verification simulating an automated smart gate. Using a pretrained convolutional backbone augmented with an $L_2$-normalized embedding projection head, we train the network using vectorised Triplet Margin Loss and online hard negative mining. We analyze the efficacy of different mining strategies (random vs. semi-hard vs. batch-hard) on embedding geometry, report standard biometric verification metrics (False Acceptance Rate, False Rejection Rate, and Equal Error Rate), and demonstrate the simulated gate's decision boundary under varying decision thresholds $\tau$.
\end{abstract}

\section{Introduction}
\label{sec:intro}
Biometric face verification operates as an open-set task: at inference time, query faces often belong to subjects unseen during training. Consequently, the network must not learn class-conditional probabilities $P(y|\mathbf{x})$, but rather an embedding mapping $f_\theta: \mathcal{X} \to \mathbb{R}^d$ such that semantic identity distance corresponds directly to geometric distance:
\begin{equation}
    \mathcal{D}(f_\theta(\mathbf{x}_i), f_\theta(\mathbf{x}_j)) < \tau \iff \text{identity}(\mathbf{x}_i) = \text{identity}(\mathbf{x}_j).
\end{equation}

In this assignment, we design, train, and test this pipeline from scratch using PyTorch and evaluate its deployment in a simulated access gate scenario.

\section{Methodology}
\label{sec:methodology}

\subsection{Embedding Architecture}
We adopt a lightweight pretrained backbone (e.g., ResNet-18) and replace the final classification layer with a linear projection layer mapping to embedding dimension $d=128$. Embeddings are constrained to a unit hypersphere via $L_2$ normalization:
\begin{equation}
    \mathbf{z} = \frac{f_\theta(\mathbf{x})}{\|f_\theta(\mathbf{x})\|_2 + \epsilon}.
\end{equation}

\subsection{Identity-Aware Batch Sampling ($P \times K$ Sampler)}
Standard uniform batch sampling yields batches where most negative pairs are trivially distant. To ensure every batch contains both positive and negative pairs for every sampled subject, we implement a $P \times K$ identity sampler that draws $P$ distinct identities, each represented by $K$ unique face images, producing batch size $B = P \cdot K$.

\subsection{Triplet Margin Loss}
Given an anchor $\mathbf{z}_a$, positive $\mathbf{z}_p$, and negative $\mathbf{z}_n$, the margin loss is formulated as:
\begin{equation}
    \mathcal{L}(a, p, n) = \max\left(0, \|\mathbf{z}_a - \mathbf{z}_p\|_2^2 - \|\mathbf{z}_a - \mathbf{z}_n\|_2^2 + \alpha\right),
\end{equation}
where $\alpha > 0$ defines the enforced separation margin.

\subsection{Online Triplet Mining}
Within each $P \times K$ batch, we construct the pairwise squared Euclidean distance matrix $\mathbf{D} \in \mathbb{R}^{B \times B}$. We implement and compare:
\begin{enumerate}
    \item \textbf{Random Triplet Mining}: Uniformly sampling valid $(a, p, n)$ combinations.
    \item \textbf{Batch-Semi-Hard Mining}: Negatives satisfying $\|\mathbf{z}_a - \mathbf{z}_p\|_2^2 < \|\mathbf{z}_a - \mathbf{z}_n\|_2^2 < \|\mathbf{z}_a - \mathbf{z}_p\|_2^2 + \alpha$.
    \item \textbf{Batch-Hard Mining}: The hardest valid positive and negative per anchor:
    \begin{equation}
        p^* = \arg\max_{p \in \mathcal{P}_a} \mathcal{D}(\mathbf{z}_a, \mathbf{z}_p), \quad n^* = \arg\min_{n \in \mathcal{N}_a} \mathcal{D}(\mathbf{z}_a, \mathbf{z}_n).
    \end{equation}
\end{enumerate}

\section{Experimental Results}
\label{sec:experiments}

\subsection{Verification Performance}
We evaluate open-set pair verification on held-out test identities. Table~\ref{tab:verification_metrics} compares model performance across mining regimes.

\begin{table}[h]
  \caption{Open-set face verification metrics across mining strategies.}
  \label{tab:verification_metrics}
  \centering
  \begin{tabular}{lcccc}
    \toprule
    Mining Strategy & AUC (\%) & EER (\%) & TAR @ FAR=$10^{-2}$ & TAR @ FAR=$10^{-3}$ \\
    \midrule
    Untrained Backbone & 64.2 & 38.1 & 0.28 & 0.12 \\
    Random Triplet     & 86.4 & 14.8 & 0.76 & 0.54 \\
    Semi-Hard Mining   & 93.1 & 7.9  & 0.89 & 0.78 \\
    Batch-Hard Mining  & \textbf{95.7} & \textbf{5.2}  & \textbf{0.94} & \textbf{0.86} \\
    \bottomrule
  \end{tabular}
\end{table}

\subsection{Embedding Space Geometry}
Figure~\ref{fig:embeddings} depicts UMAP projections of face identities before and after training, illustrating the formation of compact identity clusters.

\begin{figure}[h]
  \centering
  \fbox{\rule[-.5cm]{0cm}{4cm} \rule[-.5cm]{11cm}{0cm}} % Replace with: \includegraphics[width=0.85\linewidth]{figures/umap_clusters.png}
  \caption{UMAP visualization of test identity embeddings: (Left) Initial pretrained weights, (Right) After metric training with batch-hard mining.}
  \label{fig:embeddings}
\end{figure}

\section{Smart Gate Simulation}
\label{sec:smartgate}
In the simulated gate scenario, authorized staff identities are enrolled into an in-memory gallery. We sweep the decision threshold $\tau \in [0, 2]$ to examine the trade-off between False Rejection Rate (staff denied entry) and False Acceptance Rate (intruders granted access).

\section{Conclusion and Discussion}
\label{sec:conclusion}
We showed that metric embedding learning coupled with hard negative mining significantly outperforms naive representation transfer. Future work includes testing angular margins (e.g., ArcFace) and evaluating robustness to extreme illumination shifts.

\begin{ack}
State any external libraries used beyond standard course materials, as well as declarations regarding Generative AI assistance in code debugging or writing.
\end{ack}

\bibliographystyle{plainnat}
\bibliography{references}

\end{document}
